% ============================================================
%  ch0_front  封面、摘要、译者说明、原书页码对照表
% ============================================================

\begin{titlepage}
\centering
\vspace*{2.6cm}
{\Huge\bfseries\heiti 朗兰兹对偶 Hitchin 系统的\par}
\vspace{0.35cm}
{\Huge\bfseries\heiti 增强镜像对称\par}
\vspace{1.4cm}
{\Large Tamás Hausel\par}
\vspace{0.35cm}
{\large IST Austria\par}
\vspace{1.0cm}
{\large 单篇报告的全文汉译与译注\par}
\vspace{0.5cm}
{\large 数译\ \ 译注\par}
\vspace{1.8cm}
{\small
底本：Tamás Hausel，\emph{Enhanced mirror symmetry for Langlands dual Hitchin systems}，\\
Proc.\ Int.\ Cong.\ Math.\ 2022，Vol.~3，pp.~2228--2249，\\
DOI 10.4171/ICM2022/164，按 CC BY 4.0 授权发布。\par}
\vfill
\end{titlepage}

\section*{摘要}
\addcontentsline{toc}{section}{摘要}

本文第一部分综述 Langlands 对偶群的 Hitchin 系统之间镜像对称现象的数学结果。第二部分引入并讨论 Lagrange 子流形上 Hitchin 系统的重数代数，并考察其镜像上相应的猜想结构。

\medskip
\noindent\textbf{2020 年数学主题分类：}主类 14D21；次类 14C17、14J33、14D24、20G05、81T13。

\noindent\textbf{关键词：}Hitchin 系统，镜像对称，重数代数，几何 Langlands。

\section*{译者说明}
\addcontentsline{toc}{section}{译者说明}

\textbf{底本与校对。}\ 译文以作者为 ICM 2022 报告提供的 \LaTeX{} 源（与 arXiv:2112.09455v2 同稿）为唯一底稿，逐句对译；全部公式直接沿用源码，未作改写。原文印刷本（EMS Press，pp.~2228--2249）的断行与分页由排版决定，与本译本不同；两本内容一致，个别处印刷本略有增删，已在《疑点清单》中登记。译者另将原书页码标在每节起首处以 \origpage{2229} 这样的灰字形式给出，完整对照见下页《原书页码对照表》。

\textbf{编号系统。}\ 原文的定理、命题、推论、猜想、引理、定义、注、例共用一套计数器，按节编号；本译本沿用同一计数器与同一顺序，故编号与原文逐条相同（例如定理 3.1、定理 3.5、定理 3.12、定理 3.18 在原文与本译本中是同一定理）。公式编号同样照原样保留，凡引用“式 (3.12)”即指原文的 (3.12)。

\textbf{术语。}\ 先查本库《数语规范》（\texttt{01\_corpus\_zh/数语规范/TERMS.tsv}），有定案者遵其定案；无定案者取中文数学文献的通行译法，首次出现时以括号附英文原词；尚无通行中译者保留英文，不作音译。同一人名、同一术语在全文中只有一种写法。术语总表见附录 B。

\textbf{译注。}\ 只写三类：易错预警、含糊点拨、背景直觉。每条译注都以脚注给出，且必带 \textbf{[来源：……]}；凡属译者自身体会而无外部依据者，一律注明 \textbf{[译者按（无外部依据）]}。原文的笔误不擅自改动，照录并加注说明，同时登记在《疑点清单》中。

\textbf{文献。}\ 参考文献表照原文照录，不译标题与刊名，以便读者按原样检索。

\section*{原书页码对照表}
\addcontentsline{toc}{section}{原书页码对照表}

下表左列是原书页码（Proc.\ ICM 2022 第 3 卷的印刷页码），中列是该页起首处的内容，右列是本译本中与之对应的位置。查原文时按下表定位即可。

\begin{longtable}{p{1.5cm}p{8.4cm}p{4.3cm}}
\toprule
原书页 & 该页起首处的内容 & 本译本位置\\
\midrule
\endfirsthead
\toprule
原书页 & 该页起首处的内容 & 本译本位置\\
\midrule
\endhead
\bottomrule
\endfoot
2228 & 标题页、摘要、主题分类与关键词 & 封面、摘要\\
2229 & 一、引言：Yang--Mills 方程的二维约化，Hitchin 1987 & \S1\\
2230 & 引言末段；二、背景；2.1 镜像对称：拓扑、同调、SYZ 三个方面 & \S1 末、\S2.1\\
2231 & 2.1 末（Strominger--Yau--Zaslow）；2.2 几何 Langlands 对应 & \S2.2\\
2232 & 2.2 末（opers）；2.3 Langlands 对偶 Hitchin 系统的 SYZ 镜像对称；2.4 拓扑镜像对称 & \S2.3、\S2.4\\
2233 & 混合 Hodge 数与 $\alpha$-twisted stringy Hodge 数；猜想 2.1 & \S2.4\\
2234 & 2.5 几何 Langlands 作为增强同调镜像对称：Kapustin--Witten 的 $S$ 对偶 & \S2.5\\
2235 & 2.5 末（Wilson 算子与 Hecke 算子）；三、幂零锥尖端处的增强镜像对称 & \S2.5 末、\S3\\
2236 & 3 的动机（在零纤维上检验镜像对）；3.1 非常稳定 Higgs 丛；3.1.1 & \S3.1、\S3.1.1\\
2237 & 3.1.2 Higgs 丛的 Białynicki--Birula 分解；定理 3.1 & \S3.1.2\\
2238 & 定理 3.5；3.2 重数代数与 Lagrange 子流形上的显式 Hitchin 系统 & \S3.2\\
2239 & 定义 3.7 之后的注；定理 3.9 的准备 & \S3.2\\
2240 & 定理 3.10（重数代数即 Grassmann 上同调环）；Hecke 对应的图示 & \S3.2\\
2241 & 3.3 wobbly Lagrange 子流形上的显式 Hitchin 系统；仿射 Grassmannian & \S3.3\\
2242 & 支配权 $\mu$ 的 Hecke 对应；$\overline{\mathrm{Gr}}^{\mu}$ 的分解 & \S3.3\\
2243 & 定理 3.12 及其后的三个注（$n=2$ 的显式计算） & \S3.3\\
2244 & 3.3.1 $W^{+}_{\delta}$ 的 Lagrange 闭包；3.4 走向几何 Satake 对应的经典极限 & \S3.3.1、\S3.4\\
2245 & 式 (3.12)--(3.14)：Wilson 算子与 Hecke 算子的交换 & \S3.4\\
2246 & 猜想 3.19；注 3.20；致谢 & \S3.4 末、致谢\\
2247 & 参考文献 [1]--[12] & 参考文献\\
2248 & 参考文献 [13]--[30] & 参考文献\\
2249 & 参考文献 [31]--[45] & 参考文献\\
\end{longtable}

\newpage
\tableofcontents
\newpage
